\lysh{12628}{4}

\begin{alist}
\item Το σημείο Ζ έχει ως τετμημένη την αρνητική λύση της εξίσωσης $f(x) = 0$ η οποία γράφεται $x^2 - x - 1 = 0$ με διακρίνουσα $\Delta = (-1)^2 - 4 \cdot 1 \cdot (-1) = 1 + 4 = 5$.
Άρα $x_1 = \dfrac{-(-1)-\sqrt{5}}{2 \cdot 1} = \dfrac{1-\sqrt{5}}{2} < 0$, αφού $1 < 5 \Rightarrow \sqrt{1} < \sqrt{5} \Rightarrow 1 < \sqrt{5} \Rightarrow 1 - \sqrt{5} < 0$.
Η άλλη ρίζα είναι η $x_2 = \dfrac{-(-1)+\sqrt{5}}{2 \cdot 1} = \dfrac{1+\sqrt{5}}{2} > 0$. Ώστε $Z\left(\dfrac{1-\sqrt{5}}{2} , 0\right)$.

Το σημείο Α έχει τετμημένη μηδέν και τεταγμένη $g(0) = 3 - 0 = 3$, έτσι $A(0, 3)$.
Τα σημεία Β και Γ έχουν ως τετμημένες τις λύσεις της εξίσωσης $f(x) = g(x)$, δηλαδή της εξίσωσης $x^2 - x - 1 = 3 - x$ η οποία γράφεται $x^2 - 4 = 0$. Οπότε $x^2 = 2^2$, άρα $x = 2$ ή $x = -2$.
Οι τεταγμένες των σημείων Γ και Β θα είναι:
$f(2) = g(2) = 3 - 2 = 1$ και $f(-2) = g(-2) = 3 - (-2) = 5$ αντίστοιχα.
Ώστε $B(-2, 5)$, $\Gamma(2, 1)$ αφού το σημείο Β βρίσκεται πιο αριστερά από το Γ, άρα θα έχει μικρότερη τετμημένη.

\item Πρέπει $f(x) > g(x)$, δηλαδή $x^2 - x - 1 > 3 - x$, οπότε $x^2 > 4$, σχέση που γράφεται $x^2 > 2^2$, άρα $|x| > 2$. Ώστε $x < -2$ ή $x > 2$.
Έτσι, για $x \in (-\infty , -2) \cup (2 , +\infty)$ θα είναι $f(x) > g(x)$.

\item Προφανώς πρέπει να αποδείξουμε ότι $|f(\alpha) - (- g(\alpha))| \ge 1$ για κάθε πραγματικό αριθμό $\alpha$. Η σχέση αυτή ισοδύναμα γράφεται:
$|f(\alpha) + g(\alpha)| \ge 1$, άρα $|\alpha^2 - \alpha - 1 + 3 - \alpha| \ge 1$, δηλαδή $|\alpha^2 - 2\alpha + 2| \ge 1$. (Ι)
Παρατηρούμε ότι $\alpha^2 - 2\alpha + 2 = \alpha^2 - 2\alpha + 1 + 1 = (\alpha - 1)^2 + 1 \ge 1 > 0$, αφού για κάθε $\alpha \in \mathbb{R}$ ισχύει $(\alpha - 1)^2 \ge 0 \Leftrightarrow (\alpha - 1)^2 + 1 \ge 1$.
Γνωρίζουμε όμως ότι αν $y > 0$, τότε είναι $|y| = y$.
Έτσι η σχέση (Ι) ισοδύναμα γράφεται $|(\alpha - 1)^2 + 1| \ge 1$, δηλαδή $(\alpha - 1)^2 + 1 \ge 1$, οπότε $(\alpha - 1)^2 \ge 0$ σχέση που ισχύει για κάθε $\alpha \in \mathbb{R}$.

Εναλλακτικά, το τριώνυμο $\alpha^2 - 2\alpha + 2$ έχει αρνητική διακρίνουσα $\Delta = (-2)^2 - 4 \cdot 1 \cdot 2 = 4 - 8 = -4$, άρα το τριώνυμο γίνεται ομόσημο του συντελεστή του $\alpha^2$ δηλαδή του $1$, άρα πάντα θετικό για κάθε $\alpha \in \mathbb{R}$. Ώστε $|\alpha^2 - 2\alpha + 2| = \alpha^2 - 2\alpha + 2$. Έτσι η σχέση (Ι) γράφεται ισοδύναμα $\alpha^2 - 2\alpha + 2 \ge 1$, δηλαδή $\alpha^2 - 2\alpha + 1 \ge 0$ άρα $(\alpha - 1)^2 \ge 0$ που ισχύει για κάθε $\alpha \in \mathbb{R}$.
\end{alist}
